\section{Cuts with no crossing edge}
\label{sec:area_law_zero_boundary}

This section proves Lemma~2.2 of \cite{OpenAI2026AreaLaw} and the degenerate
cases recorded before it in Section~2 of that preprint.

\begin{theorem}[Gapped ground vectors of a sum of factor Hamiltonians]
    \label{thm:area_law_product_ground_state}
    \lean{TNLean.Lattice.exists_isGappedGroundState_factors_of_kronecker_sum, TNLean.Lattice.exists_eq_mul_of_gap, TNLean.Lattice.vonNeumannEntropy_partialTraceRight_eq_zero_of_gap}
    \leanok
    Let $H=H_A\otimes I+I\otimes H_B$ with $H_A,H_B$ Hermitian. If a unit
    vector $\omega$ satisfies $H\omega=E_0\omega$ and
    $H-E_0I\ge\Delta(I-|\omega\rangle\langle\omega|)$ with $\Delta>0$, then
    $\omega$ is a product vector, and its reduced state on the first factor has
    zero entropy.

    More precisely, there are real numbers $a,b$ and normalized factor
    vectors $u,v$ such that
    \begin{align}
        \omega & =u\otimes v,                            & E_0   & =a+b,\notag \\
        \|u\|  & =1,                                     & \|v\| & =1,\notag   \\
        H_Au   & =au,                                    & H_Bv  & =bv,\notag  \\
        H_A-aI & \ge\Delta(I-|u\rangle\langle u|),       &
        H_B-bI & \ge\Delta(I-|v\rangle\langle v|).\notag
    \end{align}
    The phase is included in the chosen factors, so the product equality
    is exact. Both factors retain the same lower gap bound $\Delta$.
    These factor gap inequalities are the spectral step in the reduction
    to connected domains in \cite[Part~I]{TNLean2026OneDimensional}.
\end{theorem}

\begin{proof}\leanok
    Let $a,b$ be the least eigenvalues of $H_A,H_B$ with unit eigenvectors
    $u,v$. Then $H\ge a+b$, so $E_0=\langle\omega,H\omega\rangle\ge a+b$, while
    $H(u\otimes v)=(H_Au)\otimes v+u\otimes(H_Bv)=(a+b)(u\otimes v)$. The gap inequality on $u\otimes v$ gives
    $\Delta(1-|\langle\omega,u\otimes v\rangle|^2)\le a+b-E_0\le0$, hence
    $|\langle\omega,u\otimes v\rangle|=1$ and $\omega$ is a multiple of
    $u\otimes v$. The reduced state is then a rank-one projection.

    The equality $|\langle\omega,u\otimes v\rangle|=1$, together with
    the two energy comparisons, also gives $E_0=a+b$.
    Write $z=\langle\omega,u\otimes v\rangle$.
    The exact relation is $\omega=\overline z\,u\otimes v$ with $|z|=1$.
    Replacing $u$ by $\overline z\,u$ preserves its norm, eigenvalue
    equation and rank-one projection and makes $\omega=u\otimes v$.
    For the isometric inclusion $B_v\xi=\xi\otimes v$, direct compression
    of the gap defect gives
    \begin{align}
         & B_v^*\bigl(H-E_0I-\Delta(I-|\omega\rangle\langle\omega|)\bigr)B_v
        \notag                                                               \\
         & \qquad=H_A-(E_0-b)I-\Delta(I-|u\rangle\langle u|)
        =H_A-aI-\Delta(I-|u\rangle\langle u|).\notag
    \end{align}
    Positivity is preserved under rectangular compression, so the last
    operator is positive semidefinite. Compressing instead along
    $B_u\eta=u\otimes\eta$ gives the second factor inequality.
    No factor is lost in $\Delta$, since both inclusions are isometries.
\end{proof}

\begin{theorem}[Supports and components]
    \label{thm:area_law_support_components}
    \lean{TNLean.Lattice.IsAdmissibleSupport.reachable, TNLean.PEPS.AreaLaw.subset_or_disjoint_of_edgeBoundary_eq_empty, TNLean.PEPS.AreaLaw.edgeBoundary_connectedComponents_eq_empty, TNLean.PEPS.AreaLaw.edgeBoundary_connectedComponent_eq_empty}
    \leanok
    \uses{def:area_law_domain,def:area_law_hamiltonian}
    Every support of finite induced-graph diameter lies in one connected
    component of the domain graph. If $\partial_\Lambda A=\varnothing$, every
    such support lies in $A$ or in $\Lambda\setminus A$. The set of sites of any
    union of connected components has no crossing edge.
\end{theorem}

\begin{proof}\leanok
    A walk inside the domain cannot pass from $A$ to its complement without
    traversing an edge of $\partial_\Lambda A$, and any two sites of a support
    are joined by such a walk. Adjacent sites lie in the same component.
\end{proof}

\begin{theorem}[A cut splitting no support]
    \label{thm:area_law_unsplit_supports}
    \lean{TNLean.PEPS.AreaLaw.regionalEntropy_eq_zero_of_forall_subset_or_disjoint}
    \leanok
    \uses{def:area_law_ground_entropy}
    Under the hypotheses of Theorem~\ref{thm:openai_uniform_area_law}, if every
    admissible support $X$ satisfies $X\subseteq A$ or $X\cap A=\varnothing$,
    then $S_\Omega(A)=0$.
\end{theorem}

\begin{proof}\leanok
    \uses{thm:area_law_product_ground_state}
    Every term of the Hamiltonian acts inside $A$ or inside its complement, so
    across the cut $H=H_A\otimes I+I\otimes H_{\Lambda\setminus A}$;
    replacing both factors by their Hermitian parts preserves this decomposition.
    Hence the ground vector is a
    product vector across the cut and its reduced state on $A$ has zero entropy.
\end{proof}

\begin{theorem}[A cut with no crossing edge]
    \label{thm:area_law_zero_boundary}
    \lean{TNLean.PEPS.AreaLaw.regionalEntropy_eq_zero_of_edgeBoundary_eq_empty}
    \leanok
    \uses{def:area_law_ground_entropy}
    Under the hypotheses of Theorem~\ref{thm:openai_uniform_area_law}, if
    $\partial_\Lambda A=\varnothing$, then $S_\Omega(A)=0$. This is
    Lemma~2.2 of \cite{OpenAI2026AreaLaw}; only the eigenvector condition and the
    full-system gap with $\Delta>0$ are used.
\end{theorem}

\begin{proof}\leanok
    \uses{thm:area_law_unsplit_supports, thm:area_law_support_components}
    By Theorem~\ref{thm:area_law_support_components}, every admissible support
    lies in $A$ or in $\Lambda\setminus A$, so
    Theorem~\ref{thm:area_law_unsplit_supports} applies.
\end{proof}

\begin{theorem}[Degenerate cases]
    \label{thm:area_law_degenerate_cases}
    \lean{TNLean.PEPS.AreaLaw.regionalEntropy_eq_zero_of_range_eq_zero, TNLean.PEPS.AreaLaw.regionalEntropy_eq_zero_of_q_eq_one, TNLean.PEPS.AreaLaw.regionalEntropy_eq_zero_of_domain_eq_empty, TNLean.PEPS.AreaLaw.regionalEntropy_connectedComponents_eq_zero}
    \leanok
    \uses{def:area_law_ground_entropy}
    For range $R=0$ every regional entropy of the gapped ground vector vanishes.
    For $q=1$ or an empty domain every regional entropy of every unit vector
    vanishes. Every union of connected components, and in particular the empty
    region and the full domain, has zero regional entropy for the gapped ground
    vector.
\end{theorem}

\begin{proof}\leanok
    \uses{thm:area_law_unsplit_supports, thm:area_law_zero_boundary, thm:area_law_support_components}
    For $R=0$ every support is a single site and so is never split by a cut,
    so Theorem~\ref{thm:area_law_unsplit_supports} applies.
    For $q=1$ or $\Lambda=\varnothing$ the regional Hilbert space is
    one-dimensional, and a unit-trace state on it has zero entropy. A union of
    connected components has no crossing edge by
    Theorem~\ref{thm:area_law_support_components}, so
    Theorem~\ref{thm:area_law_zero_boundary} applies.
\end{proof}
